% DRAFT deviation table (Efe): changes since the approved proposal, each with
% its reason. Needs \usepackage{booktabs,tabularx} in the preamble.
\begin{table*}[t]
\centering
\caption{Deviations from the approved proposal and their justification.}
\label{tab:deviations}
\footnotesize
\begin{tabularx}{\textwidth}{@{}p{2.6cm}XXp{2.2cm}@{}}
\toprule
\textbf{Aspect} & \textbf{Proposal} & \textbf{Now, and why} & \textbf{Trigger} \\
\midrule
Meaning of $\sigma$ &
$\sigma$ predicted by the NLL head expected to track listener disagreement (H2). &
$\sigma$ is defined as the predicted error of the model on the \emph{mean} rating
(model error plus averaging noise $\approx s^2/n$), not the spread of individual
listeners. $\rho(\sigma,s)$ is kept only as a descriptive analysis. H2 is split
into H2a (calibration, $\sigma$ vs.\ actual error), H2b (how measurable $s$ is)
and H2c (a head trained on individual ratings, after the report). &
Instructor feedback; Watcharasupat \emph{et al.}~\cite{watcharasupat2025} \\
Individual ratings &
Not considered. &
The data audit found per-rater static ratings (10--20 raters per clip, 187 raters),
which enables reliability ceilings now and explicit modelling of the rating
distribution (H2c) later. &
Instructor feedback; data audit \\
Split protocol &
One artist-disjoint 70/15/15 split, three seeds, bootstrap over test songs. &
Artist-grouped, quadrant-stratified 10-fold CV; fold $k{+}1$ is validation.
Every clip is tested exactly once, so uncertainty analyses use all 1{,}744
out-of-fold predictions instead of $\approx$262, and models are compared by
paired per-fold differences. &
Small dataset; \cite{bouthillier2021} \\
Success thresholds &
Fixed $R^2$ targets (e.g.\ C1: $R^2_V\geq0.30$, $R^2_A\geq0.60$; C2: $R^2_V\geq0.45$, $R^2_A\geq0.68$). &
C1 becomes a sanity check reported with its spread, not a pass/fail gate. The
primary H1 criterion is relative: the pretrained encoder must beat openSMILE
\emph{with the same head} on both dimensions. &
Instructor feedback (thresholds optimistic under artist-disjoint CV) \\
Same-head comparison &
SVR on openSMILE vs.\ MLP on MERT. &
Ridge and MLP heads are trained on both openSMILE and MERT so that the
representation effect is not confounded with the head; SVR is reported as an
additional baseline. &
Methodological \\
openSMILE excerpt &
DEAM features used as supplied. &
For the 744 clips of 2013 the supplied feature files cover the \emph{full song}
(up to 585\,s) while audio and ratings cover a 45\,s excerpt. Features are now
cropped to the excerpt using the segment start in the 2013 metadata (offsets
verified against the audio energy). &
Found during baseline work \\
2013 static labels &
Labels used as supplied. &
For the 744 clips of 2013 the static valence and arousal columns are swapped
relative to the dynamic annotations ($r\approx0.95$ crossed vs.\ $\approx0.57$
matched; 2014 is consistent) and to loudness. They are swapped back. &
Found during baseline work \\
2015 full songs &
Secondary out-of-distribution test set (58 songs). &
Dropped: different annotation protocol and only 5 raters per song. &
Scope \\
Scope of the report &
PANNs, fine-tuning and ensembles in weeks 3--5. &
PANNs, random-init MERT (null model), layer probe and LP-FT moved after the
report, following the instructor's ordering of priorities. &
Instructor feedback \\
\bottomrule
\end{tabularx}
\end{table*}
